
\begin{obeassessment}
\textbf{Kuis Bab 5}
1. Apa perbedaan utama antara \textit{true random} dan \textit{pseudo-random} dalam konteks \textit{keystream generator}?
2. Mengapa operasi XOR dipilih sebagai operasi utama dalam \textit{stream cipher}?
3. Jelaskan prinsip kerja \textit{majority clocking} pada algoritma A5/1.
4. Hitunglah cipherteks dari plainteks \texttt{11010110} yang di-XOR dengan kunci \texttt{01101100}.
5. Sebutkan tiga syarat yang harus dipenuhi agar keystream generator dapat dipakai berulang kali dengan kunci yang sama untuk pesan berbeda.
6. Jelaskan mengapa RC4 yang menerima kunci sampai 2048 bit tetap dapat dipecahkan, sedangkan panjang kunci justru dianggap sebagai ukuran keamanan pada cipher lain.
7. Apa fungsi pemanasan 1152 detak pada Trivium, dan apa yang terjadi bila tahap itu dihilangkan?
\end{obeassessment}
\begin{competencychecklist}
\begin{tabularx}{\linewidth}{|X|c|}
\hline
Indikator Kompetensi & Tercapai \\
\hline
Saya dapat menjelaskan perbedaan kriptografi klasik dan modern & [ ] \\
\hline
Saya dapat mengonversi data antara biner, byte, dan heksadesimal & [ ] \\
\hline
Saya dapat melakukan operasi bitwise XOR & [ ] \\
\hline
Saya dapat menjelaskan mekanisme kerja \textit{stream cipher} & [ ] \\
\hline
Saya dapat menjelaskan peran \textit{key}, \textit{nonce}, dan \textit{counter} pada keystream generator & [ ] \\
\hline
Saya dapat menghitung keystream LFSR beserta periodenya & [ ] \\
\hline
Saya dapat menjelaskan mengapa LFSR tunggal tidak aman & [ ] \\
\hline
Saya dapat menjelaskan tahap KSA dan PRGA pada RC4 serta kelemahannya & [ ] \\
\hline
Saya dapat menjelaskan kaidah mayoritas dan pola inisialisasi A5/1 & [ ] \\
\hline
Saya dapat menjelaskan struktur NLFSR dan pemanasan pada Trivium & [ ] \\
\hline
Saya dapat menjelaskan operasi AXR dan \textit{quarter round} ChaCha20 & [ ] \\
\hline
Saya dapat membedakan karakteristik RC4, A5/1, dan ChaCha20 & [ ] \\
\hline
Saya dapat menganalisis keuntungan \textit{stream cipher} untuk data \textit{real-time} & [ ] \\
\hline
Saya dapat menunjukkan bahaya pemakaian ulang keystream melalui perhitungan XOR & [ ] \\
\hline
\end{tabularx}
\end{competencychecklist}
